\section{Conclusion}
\label{sec:conclusion}
In this work, we investigate the problem of vulnerability type identification (VTI) after vulnerability detection. The problem is modeled as the multi-label classification task. Particularly, each detected vulnerable function is considered as a document, and the set of vulnerability types of the function could be considered as the label set of the corresponding document.
%
This NLP task has been effectively addressed by pre-training, then fine-tuning the framework with deep pre-trained embedding models. 
%
The existing studies show that the deep pre-trained embedding models specialized for code have also shown their effectiveness for many classification tasks in software engineering.
%
In this paper, we experimentally evaluate the performance of the well-known and advanced pre-trained models for VTI on a large set of vulnerabilities in various types. Surprisingly, their performance is not much better than the VTI performance of the traditional baseline classification model with an old-fashioned bag-of-word TF-IDF. Meanwhile, these neural network approaches cost much more time and require GPU. We also introduce a lightweight independent component to enhance the predictions of the baseline approach. Our idea is that the types of vulnerabilities could strongly correlate to certain code tokens (distinguishing tokens) in several crucial parts of programs. The distinguishing tokens for a \vt are statistically identified based on their prevalence in the type versus the others. Our results show that the baseline approach enhanced by our component can outperform the state-of-the-art deep pre-trained methods while retaining very high efficiency. Furthermore, the proposed technique could also improve the neural network approaches by up to 92.8\% in macro-average F1.